\documentclass[tikz,border=6pt]{standalone}
\usepackage{mathdiagrams}
\begin{document}
\begin{tikzpicture}[md]
  % Triangle ABC with exterior angle at C (extend BC to D)
  % Given: ∠A = 25°, ∠B = (x+15)°, exterior ∠ACD = (3x−10)°
  \coordinate (A) at (0,0);
  \coordinate (B) at (4.6,0);
  \coordinate (C) at (1.55,2.35);
  \coordinate (D) at (6.2,0); % B extended past C? No — extend side BC beyond C
  % Actually exterior at C: extend BC past C. Easier: extend AC? Lesson says exterior opposite to the two remote interiors at A and B is at C.
  % Exterior angle ACD: D on line BC beyond C.
  \draw[md shape] (A) -- (B) -- (C) -- cycle;
  % extend BC through C to D
  \draw[md dashed] (B) -- ($(B)!1.55!(C)$) coordinate (D);
  \draw[md ray] (C) -- (D);

  \node[md point] at (A) {};
  \node[md point] at (B) {};
  \node[md point] at (C) {};
  \node[md label] at (A) [below left=1pt] {$A$};
  \node[md label] at (B) [below right=1pt] {$B$};
  \node[md label] at (C) [above right=1pt] {$C$};
  \node[md label] at (D) [right=2pt] {$D$};

  % interior at A = 25°
  \pic[md angle arc, angle radius=0.65cm, "$25^\circ$", angle eccentricity=1.35] {angle=B--A--C};
  % interior at B = (x+15)°
  \pic[md angle arc, angle radius=0.65cm, "$(x+15)^\circ$", angle eccentricity=1.4] {angle=C--B--A};
  % exterior at C between CA and CD
  \pic[md highlight, angle radius=0.75cm, "$(3x-10)^\circ$", angle eccentricity=1.35] {angle=D--C--A};
  % small mark for interior at C
  \pic[md thin, color=mdMuted, angle radius=0.4cm, "", angle eccentricity=1] {angle=A--C--B};
\end{tikzpicture}
\end{document}
